% Glossário do capítulo — no fim do capítulo, como consulta (v3 §5)
\section*{Glossário do capítulo}
\addcontentsline{toc}{section}{Glossário do capítulo}
\markright{Glossário}
{\small Termos e símbolos usados no capítulo, com o significado em uma linha e a unidade onde são construídos.}

\subsection*{Termos}
\begin{glossario}
\item[Proposição (ou sentença)] Oração declarativa que tem exatamente um valor lógico: verdadeiro (V) ou falso (F). \hfill\emph{Unid.~1}
\item[Valor lógico] O V ou o F de uma proposição. \hfill\emph{Unid.~1}
\item[Negação] Operação que forma, a partir de $p$, a proposição ${\sim}p$, de valor oposto. \hfill\emph{Unid.~1}
\item[Tabela-verdade] Tabela que lista todos os casos de valor das proposições simples e o valor resultante. \hfill\emph{Unid.~1}
\item[Sentença aberta] Frase com variável sem valor atribuído; ainda não é proposição. \hfill\emph{Unid.~1 e 5}
\item[Princípio da não contradição] Uma proposição não é verdadeira e falsa ao mesmo tempo. \hfill\emph{Unid.~1}
\item[Princípio do terceiro excluído] \enquote{$p$ ou não $p$} é sempre verdadeira. \hfill\emph{Unid.~1}
\item[Bivalência] A tese de que toda proposição tem um, e só um, dos dois valores. \hfill\emph{Unid.~1}
\item[Contraditórias] Duas frases que nunca têm o mesmo valor: uma é a negação da outra (\enquote{todo $S$ é $P$} × \enquote{algum $S$ não é $P$}). \hfill\emph{Unid.~1 e 6}
\item[Contrárias] Duas frases que não podem ser ambas verdadeiras, mas podem ser ambas falsas (\enquote{todo $S$ é $P$} × \enquote{nenhum $S$ é $P$}). \hfill\emph{Unid.~1 e 6}
\item[Paradoxo] Argumento que, de premissas aparentemente aceitáveis, chega a uma contradição. \hfill\emph{Unid.~1}
\item[Concordância negativa] Fenômeno do português em que várias palavras negativas formam uma única negação (\enquote{não vi ninguém}). \hfill\emph{Unid.~1}
\item[Conectivo] Símbolo que forma uma proposição composta a partir de outras ($\wedge$, $\vee$, $\rightarrow$, $\leftrightarrow$). \hfill\emph{Unid.~2 e 3}
\item[Proposição composta] Proposição formada por outras mediante conectivos ou negação. \hfill\emph{Unid.~2}
\item[Conjunção] A proposição $p\wedge q$ (\enquote{$p$ e $q$}): verdadeira só quando $p$ e $q$ são ambas verdadeiras. \hfill\emph{Unid.~2}
\item[Disjunção (inclusiva)] A proposição $p\vee q$ (\enquote{$p$ ou $q$}): falsa só quando $p$ e $q$ são ambas falsas. \hfill\emph{Unid.~2}
\item[Disjunção exclusiva] \enquote{Ou $p$ ou $q$, mas não ambos}: verdadeira quando $p$ e $q$ têm valores diferentes. Não é conectivo do FME. \hfill\emph{Unid.~2}
\item[Verofuncionalidade] Propriedade de um conectivo cujo resultado depende só dos valores das partes, não do seu assunto. \hfill\emph{Unid.~2}
\item[Condicional] A proposição $p\rightarrow q$ (\enquote{se $p$, então $q$}): falsa só quando $p$ é V e $q$ é F. \hfill\emph{Unid.~3}
\item[Antecedente e consequente] Em $p\rightarrow q$, $p$ é o antecedente (hipótese) e $q$ o consequente (tese). \hfill\emph{Unid.~3}
\item[Condição suficiente / necessária] Em $p\rightarrow q$ verdadeiro, $p$ é suficiente para $q$ e $q$ é necessária para $p$. \hfill\emph{Unid.~3}
\item[Bicondicional] A proposição $p\leftrightarrow q$ (\enquote{$p$ se, e somente se, $q$}): verdadeira quando $p$ e $q$ têm o mesmo valor. \hfill\emph{Unid.~3}
\item[Recíproca, inversa, contrapositiva] De $p\rightarrow q$: $q\rightarrow p$, ${\sim}p\rightarrow{\sim}q$ e ${\sim}q\rightarrow{\sim}p$. Só a contrapositiva equivale ao original. \hfill\emph{Unid.~3}
\item[Implicação material] Nome dado por Russell ao condicional de Fílon, o da tabela de $\rightarrow$. \hfill\emph{Unid.~3}
\item[Tautologia] Proposição composta verdadeira em todas as linhas da tabela (logicamente verdadeira). \hfill\emph{Unid.~4}
\item[Contradição] Proposição composta falsa em todas as linhas (logicamente falsa). \hfill\emph{Unid.~4}
\item[Contingência] Proposição composta com V e F na última coluna. Termo da literatura; o FME não o usa. \hfill\emph{Unid.~4}
\item[Implicação ($\Rightarrow$)] Relação: $p\Rightarrow q$ quando não há linha com $p$ V e $q$ F, isto é, quando $p\rightarrow q$ é tautologia. \hfill\emph{Unid.~4}
\item[Equivalência ($\Leftrightarrow$)] Relação: $p\Leftrightarrow q$ quando $p$ e $q$ têm a mesma tabela, isto é, quando $p\leftrightarrow q$ é tautologia. \hfill\emph{Unid.~4}
\item[Argumento válido] Argumento em que (premissas) $\Rightarrow$ (conclusão). \hfill\emph{Unid.~4}
\item[Quantificador universal ($\forall$)] \enquote{Para todo}: $(\forall x)P(x)$ é V quando $P(x)$ vale para todos os valores do domínio. \hfill\emph{Unid.~5}
\item[Quantificador existencial ($\exists$)] \enquote{Existe (pelo menos um)}: $(\exists x)P(x)$ é V quando $P(x)$ vale para algum valor. \hfill\emph{Unid.~5}
\item[Existe um único ($\exists|$)] Existe exatamente um valor com $P(x)$; também se escreve $\exists!$. \hfill\emph{Unid.~5}
\item[Domínio (ou universo)] Conjunto em que a variável assume valores; o valor de uma quantificada depende dele. \hfill\emph{Unid.~5}
\item[Contraexemplo] Caso particular que torna falsa uma afirmação do tipo \enquote{para todo}. \hfill\emph{Unid.~3 e 5}
\item[Função proposicional] Outro nome da sentença aberta: a cada valor da variável associa uma proposição. \hfill\emph{Unid.~5}
\item[Quadrado das oposições] Diagrama das relações entre as formas A, E, I, O (todo, nenhum, algum, algum\ldots\ não). \hfill\emph{Unid.~6}
\item[Leis de De Morgan] ${\sim}(p\wedge q)\Leftrightarrow{\sim}p\vee{\sim}q$ e ${\sim}(p\vee q)\Leftrightarrow{\sim}p\wedge{\sim}q$. \hfill\emph{Unid.~4 e 6}
\item[Redução ao absurdo] Prova que supõe a negação do que se quer provar e chega a uma contradição. \hfill\emph{Unid.~6}
\item[Álgebra de Boole] As leis de equivalência dos conectivos (comutativas, associativas, distributivas, absorção, De Morgan\ldots). \hfill\emph{Unid.~4}
\end{glossario}

\subsection*{Símbolos}
\begin{center}\small
\begin{tabular}{clll}\toprule
\sffamily\bfseries Símbolo & \sffamily\bfseries Leitura & \sffamily\bfseries Exemplo & \sffamily\bfseries Unid.\\\midrule
${\sim}$ & não & ${\sim}p$ & 1\\\fino
$\neg$ & não (notação alternativa) & $\neg p$ & 1\\\fino
$\in$,\ $\notin$ & pertence, não pertence & $2\in\mathbb{Z}$,\ $\sqrt2\notin\mathbb{Q}$ & 1\\\fino
$\mid$,\ $\nmid$ & divide, não divide & $3\mid 12$,\ $3\nmid 11$ & 1\\\fino
$\wedge$ & e (conjunção) & $p\wedge q$ & 2\\\fino
$\vee$ & ou (disjunção inclusiva) & $p\vee q$ & 2\\\fino
$\veebar$ & ou exclusivo (auxiliar; não é do FME) & $p\veebar q$ & 2\\\fino
$\cap$,\ $\cup$ & interseção, união & $x\in A\cap B$ & 2 e Cap.~II\\\fino
$\rightarrow$ & se\ldots, então\ldots & $p\rightarrow q$ & 3\\\fino
$\leftrightarrow$ & \ldots se, e somente se,\ldots & $p\leftrightarrow q$ & 3\\\fino
$\Rightarrow$ & implica (relação) & $p\wedge q\Rightarrow p$ & 4\\\fino
$\Leftrightarrow$ & é equivalente a (relação) & $p\rightarrow q\Leftrightarrow{\sim}p\vee q$ & 4\\\fino
$\forall$ & para todo, qualquer que seja & $(\forall x)(x^2\ge0)$ & 5\\\fino
$\exists$ & existe, existe pelo menos um & $(\exists x)(x+1=7)$ & 5\\\fino
$\exists|$\ \ ($\exists!$) & existe um único & $(\exists|\,x)(x+1=7)$ & 5\\
\bottomrule
\end{tabular}
\end{center}
\clearpage
